\documentclass[twocolumn]{aastex631}
\begin{document}
\title{A residual reionization ionizing-photon-budget shortfall for star-forming galaxies at z\~{}8 (highC systematic corner)}
\author{NebulaMind Lab (autonomous pipeline)}
\affiliation{NebulaMind Lab, \url{https://lab.nebulamind.net}}
\begin{abstract}
Reionization ionizing-photon-budget reconciliation at z\~{}8: star-forming galaxies require f\_esc=0.261 (+0.244/-0.129) to close the budget (Madau-Dickinson SFRD, log xi\_ion=25.5+/-0.15, clumping C=2-5, JWST-SFRD tail), versus indirect-proxy-inferred f\_esc=0.062 (+0.108/-0.039) from LzLCS O32/beta calibrations. Median delta(required-inferred)=+0.179 dex-frac (16-84\%: +0.027 to +0.424); 88\% of the systematic MC shows a shortfall. Conclusion: a genuine SHORTFALL remains, and the sign holds under both O32 and beta calibrations. The result is bounded by the xi\_ion x clumping x proxy-calibration systematic, not by statistics. Generated autonomously from public data (jwst) via the NebulaMind Lab runner: a bounded, reproducible, descriptive study.
\end{abstract}
\section{Introduction}
The reionization ionizing-photon-budget crisis has been a topic of interest in recent years, with studies suggesting that star-forming galaxies may not produce enough photons to drive reionization [Muñoz2024]. This issue is further complicated by the need for accurate measurements of escape fractions (f\_esc) and other factors contributing to the photon budget. Previous research has highlighted the importance of considering various systematics in these calculations, such as clumping and ionizing efficiency [Davies2021].
\section{Data and method}
In this study, we address the reionization-photon-budget crisis by performing a literature-anchored budget calculation. We use the cosmic star formation rate density (SFRD) from Madau \& Dickinson (2014), along with published values for xi\_ion and O32/beta f\_esc proxy calibrations [Chisholm+22, Flury+22; Simmonds+24]. Our method focuses on reconciling systematics over published literature values without relying on new observational or catalog data.
\section{Result}
Our analysis reveals that star-forming galaxies at z\~{}8 require an escape fraction of f\_esc=0.261 (+0.244/-0.129) to close the ionizing-photon-budget, assuming a Madau-Dickinson SFRD, log xi\_ion=25.5+/-0.15, and clumping factor C=2-5. However, indirect-proxy-inferred f\_esc from LzLCS O32/beta calibrations yields a significantly lower value of 0.062 (+0.108/-0.039). This discrepancy results in a median shortfall of +0.179 dex-frac (16-84\%: +0.027 to +0.424), with 88\% of systematic Monte Carlo simulations showing a deficit.
\begin{figure}\centering\includegraphics[width=\columnwidth]{result.png}\caption{A residual reionization ionizing-photon-budget shortfall for star-forming galaxies at z\~{}8 (highC systematic corner)}\end{figure}
\section{Caveats}
It is essential to acknowledge the limitations of our approach, which relies on automated, single-selection, and uncalibrated measurements. These constraints may introduce biases and uncertainties in our calculations, as they do not account for variations in galaxy properties or observational errors. Additionally, our reliance on published literature values means that any systematic errors present in those studies could propagate into our results. Therefore, further research is needed to refine these estimates and better understand the reionization-photon-budget crisis.
\section*{References}
\footnotesize
[Muoz2024] Muñoz et al. (2024). Reionization after JWST: a photon budget crisis?. bibcode 2024MNRAS.535L..37M \par [Davies2021] Davies et al. (2021). The Predicament of Absorption-dominated Reionization: Increased Demands on Ionizing Sources. bibcode 2021ApJ...918L..35D \par [Park2022] Park et al. (2022). Calibrating excursion set reionization models to approximately conserve ionizing photons. bibcode 2022MNRAS.517..192P \par [Duncan2015] Duncan et al. (2015). Powering reionization: assessing the galaxy ionizing photon budget at z < 10. bibcode 2015MNRAS.451.2030D \par [Madau2017] Madau (2017). Cosmic Reionization after Planck and before JWST: An Analytic Approach. bibcode 2017ApJ...851...50M

\end{document}
